\section{Two coordinate forms and two finite realization methods}
\label{c:rational-dictionary}

The rational construction already gives a finite cover with a strict
area saving.  We now express the same geometry in fully orthonormal
coordinates and compare two other ways to obtain its finite cylinders.
The coordinate change tracks the slope, quartic padding, and
perpendicular area together.  In these coordinates the logarithmic-cell
theorem gives bounded open bases; in the centered coordinates the
product-cell theorem does so under a weaker smoothness assumption.
All three methods approximate the same weighted integral cost to
arbitrary additive accuracy, but their hypotheses and base shapes differ.

Put $c_*=1/\sqrt2$ and
\begin{equation}\label{c:uncentered-tetrahedron}
 K_0=\{(x_0,y_0,z_0):0\le z_0\le2,
          |y_0|\le c_*z_0,\ |x_0|\le c_*(2-z_0)\}.
\end{equation}
It is a regular tetrahedron of edge $2\sqrt2$. For a fixed
$d_0>0$, define on $s>-d_0$
\begin{align}
 g(s)&=\frac{1+d_0}{s+d_0},\notag\\
 v_1(u,s)&=c_*^2s g(s)+\tfrac12(c_*^2s^2-u^2)g'(s),
 &v_2(u,s)&=u g(s).\label{c:uncentered-field}
\end{align}
This is the stream-function field associated with
$g(s)(u^2-c_*^2s^2)/2$. In particular it is tangent to
$u=\pm c_*s$. Let
\begin{align}
 A_0&=c_*^2\left(1-\frac1{2(1+d_0)}\right),&
 B_0&=\frac1{2(1+d_0)},\notag\\
 a_0(u)&=A_0u^2+B_0u^4=u^2v_1(u,1),\label{c:uncentered-polynomial}\\
 U_k&=\{(u,s):0\le s\le3/2, |u|\le c_*s,
                   s\le1+k^2(a_0(u)+\eta_0)\},\notag
\end{align}
where $\eta_0>0$ will be fixed before the tilt $k>0$.
For $|r|\le2c_*$, write $F_r(p)=p+rk v(p)$. The two swept sets are
\begin{align}
 E_1&=\{(x_0,y_0,z_0):|x_0|\le2c_*, (y_0,z_0)\in F_{x_0}(U_k)\},\notag\\
 E_2&=\{(x_0,y_0,z_0):|y_0|\le2c_*, (x_0,2-z_0)\in F_{-y_0}(U_k)\}.
 \label{c:uncentered-sweeps}
\end{align}

\begin{proposition}\label{c:rational-coordinate-identity}
Let $\ell=1/(2\sqrt2)$ and define the Euclidean similarity
\[
 \Phi(x_0,y_0,z_0)=(-\ell x_0,\ell y_0,\ell(1-z_0)).
\]
Then $\Phi K_0=K_1$, and the two sets in
\eqref{c:uncentered-sweeps} are carried exactly to the two sets in
\eqref{c:rat-sweeps} with
\begin{equation}\label{c:rational-parameter-dictionary}
 d=d_0/2,\qquad \epsilon=4(1+d_0)k,
 \qquad q=\eta_0 k^2/2.
\end{equation}
In particular $q=\eta\epsilon^2$ where
$\eta=\eta_0/[32(1+d_0)^2]$. All perpendicular areas are multiplied
by $1/8$.
\end{proposition}
\begin{proof}
Write $(Y,T)=(\ell u,(1-s)/2)$, so $r=s/2$.
Substitution into \eqref{c:rat-field} gives
\begin{equation}\label{c:rational-field-dictionary}
 v_1(u,s)=-4(1+d_0)a(Y,T),\qquad
 v_2(u,s)=2\sqrt2(1+d_0)b(Y,T).
\end{equation}
At $s=1$, these identities and
$\alpha=(1+d_0)^{-1}$ give
\[
 P(\ell u)=\frac{a_0(u)}{32(1+d_0)^2}.
\]
Consequently the last inequality defining $U_k$ is exactly
$T\ge-\epsilon^2P(Y)-q$, and $s\le3/2$ is $T\ge-1/4$.
The wedge inequalities correspond as well. For the first family,
$x=-\ell x_0$ converts its two coordinate equations into
$y=Y+\epsilon x a(Y,T)$ and $t=T+\epsilon x b(Y,T)$.
For the second, take $(X,R)=(-\ell u,(1-s)/2)$ and use that $a$
is even and $b$ is odd in the first coordinate. This gives precisely
the second equation in \eqref{c:rat-sweeps}. The bounds on the free
coordinate become $|x|,|y|\le1/2$.
Finally, $\Phi$ is an isometry followed by dilation by $\ell$;
it multiplies every perpendicular area by $\ell^2=1/8$.
\end{proof}

The monotone thresholds in Section~\ref{c:rat-threshold} therefore
give, in these coordinates, the sufficient condition
\begin{equation}\label{c:uncentered-threshold}
 z_*\le 1+kr u+k^2\bigl(a_0(u)-r^2v_1(u,1)\bigr)
       \quad\Longrightarrow\quad (u,z_*)\in F_r(U_k),
\end{equation}
provided $0\le z_*\le2$ and $|u|\le c_*z_*$.
The two right-hand thresholds sum to
$2+k^2 B_0(x_0^2-y_0^2)^2\ge2$.
Thus the exact transformation preserves both the bounded contraction
argument and the nonnegative overlap surplus, including the wedge
walls.

\subsection{Logarithmic-cell realization of the fixed-padding cover}

For the field \eqref{c:uncentered-field}, direct differentiation gives
\begin{equation}\label{c:uncentered-spectrum}
 \tr Dv=0,\qquad
 \det Dv=-c_*^2\left(\frac{(1+d_0)d_0}{(s+d_0)^2}\right)^2<0.
\end{equation}
The determinant simplification uses $gg''/2=(g')^2$.
For fixed $d_0$, the field is $C^2$ on an open rectangle containing
$U_k$ for every sufficiently small $k$. The boundary of $U_k$ has
area zero, since it lies in finitely many lines and a polynomial graph.
The strict spectral condition
$2c_* k\sqrt{-\det Dv}<1$ also holds there for small $k$.
These are exactly the hypotheses of \BLogarithmic{}.
It supplies finite covers of each set in \eqref{c:uncentered-sweeps}
with bounded open measurable bases, at cost arbitrarily close from above to
\begin{equation}\label{c:logarithmic-cost}
 J_k=\int_{U_k}(1+k^2|v(u,s)|^2)^{-1/2}\,du\,ds.
\end{equation}
The second application uses $-v$ and an isometry.

The exact cost comparison can also be checked in these coordinates.
Put $U_0=\{0\le s\le1, |u|\le c_*s\}$, of area $c_*$.
The same strip calculation as in Section~\ref{c:rat-area} gives
\begin{equation}\label{c:uncentered-cost}
 J_k=c_*+k^2\left\{
 \int_{-c_*}^{c_*}(a_0(u)+\eta_0)\,du
 -\tfrac12\int_{U_0}|v|^2\right\}+O(k^4).
\end{equation}
For fixed $d_0,\eta_0$, the added horizontal and vertical extensions
are $O(k^2)$; their product is the fourth-order boundary remainder.
Letting $d_0\downarrow0$ only in the coefficient integrals, and setting
$u=sb$, gives
\begin{align*}
 \tfrac12\int_{U_0}|v|^2&\longrightarrow
 \frac14\int_{-c_*}^{c_*}
       \left[\frac{(c_*^2+b^2)^2}{4}+b^2\right]db
       =\frac{c_*^3}{6}+\frac{7c_*^5}{30},\\
 \int_{-c_*}^{c_*}a_0(u)\,du&\longrightarrow\frac{8c_*^5}{15}.
\end{align*}
Bounded convergence is valid because $|u|\le c_*s$ bounds the fields
uniformly for $0<d_0<1$. The first limit exceeds the second by
$c_*^3/60>0$. Choose $d_0>0$ such that the actual difference
\[
 \Delta(d_0)=\tfrac12\int_{U_0}|v|^2
                  -\int_{-c_*}^{c_*}a_0(u)\,du
\]
is positive. Then choose $0<\eta_0<\Delta(d_0)/(2c_*)$,
and finally $k$ small. Formula~\eqref{c:uncentered-cost}
gives $J_k<c_*$. Only then choose the two finite errors with sum
less than $2(c_*-J_k)$. The resulting finite open-base cover has
cost below $2c_*=A_{\min}(K_0)/2$.

The $C^2$ rectangle and strict spectral inequality give the open-base
conclusion directly in the orthonormal $(u,s)$ coordinates.
The scaling identity gives
$C_d(\epsilon,q)=J_k/8$, in agreement with
Proposition~\ref{c:rational-coordinate-identity}.

\subsection{Product-cell realization of the centered cover}

For $q=\epsilon^{5/2}$, use the original centered label set
$D_{\epsilon,q}$ and the two slopes $V=\pm\epsilon W_d$.
The set is compact with area-zero boundary, $V$ is $C^1$ on an open
neighborhood, and $DV$ has eigenvalues
$\pm\epsilon d/[4(r+d)^2]$. With $M=1/2$ and $h=1/\sqrt2$,
the strict condition $0<M\lambda<1$ follows after fixing $d>0$
and taking $\epsilon$ sufficiently small. Therefore \BProduct{}
gives each family a finite cover with bounded open measurable bases
and cost at most $C_d(\epsilon,\epsilon^{5/2})+e$.

The product-cell proof chooses its padding, product threshold, and
local scale in that order, exhausts the interior by disjoint cells,
and covers the remaining null set before taking a finite open subcover.
These approximation choices occur after the present geometric
parameters are fixed. By \eqref{c:rat-cost-expansion}, choose $d$
with $I_d>J_d$, then $\epsilon$ so that $C_d<h/8$, and finally
the two errors with sum below $h/4-2C_d$. This gives the asserted
open-base cover without altering either the fractional padding or
the weighted physical metric. The compact parallelogram construction
in Section~\ref{c:rat-area} remains an independent finite realization
of the same swept sets.
